\chapter{Vladimir Voevodsky}\label{ch2002vladimirvoevodsky:vladimir-voevodsky-fields-medal-2002-a-comprehensive-report-on-his-work-and-impact}

\begin{center}\rule{0.5\linewidth}{0.5pt}\end{center}

\textbf{Prepared:} July 2026 \textbf{Subject:} Vladimir Voevodsky (June 4, 1966 -- September 30, 2017, Moscow, Russia) \textbf{Position:} Professor, Institute for Advanced Study, Princeton (2002--2017) \textbf{Doctoral Advisor:} David Kazhdan (Harvard University) \textbf{Fields Medal Citation:} \emph{``For his development of the homotopy theory of schemes and the proof of the Milnor conjecture.''}

\begin{center}\rule{0.5\linewidth}{0.5pt}\end{center}

\begin{prereqs} Algebraic geometry (motives, \u00e9tale cohomology, Nisnevich topology), homotopy theory (model categories, spectra), basic category theory, algebraic K-theory. \end{prereqs}

\begin{objectives} Construct the motivic cohomology groups $H^{p,q}(X, \mathbb{Z})$ and relate them to Milnor K-theory; prove the Milnor conjecture $K^M_n(F)/2 \cong H^n(F, \mathbb{Z}/2\mathbb{Z})$; explain the $\mathbb{A}^1$-homotopy category and the univalence axiom in homotopy type theory. \end{objectives}

\section{Problem History}

Before Voevodsky's homotopy theory \cite{voevodsky2002}, algebraic geometry lacked a satisfactory theory of \emph{motivic cohomology}---a universal cohomology theory for algebraic varieties that would unify the various Weil cohomology theories (singular, \u00e9tale, de Rham). Grothendieck conjectured the existence of such a theory in the 1960s, but constructing it remained an open problem for decades. The category of motives, his proposed universal target for cohomology functors, was defined only indirectly through the category of pure motives built from algebraic cycles.

Independently, John Milnor formulated his conjecture in 1970, relating the Milnor K-theory $K^M_*(F)$ of a field to its Galois cohomology $H^*(F, \mathbb{Z}/2\mathbb{Z})$. The Milnor K-theory is generated by symbols $\{a_1, \ldots, a_n\}$ with $a_i \in F^\times$, subject to the Steinberg relation $\{a, 1-a\} = 0$. The conjecture posits an isomorphism:
\[
K^M_n(F)/2 \cong H^n(F, \mathbb{Z}/2\mathbb{Z}).
\]
By the 1990s, the conjecture had been verified for fields of low transcendence degree over $\mathbb{Q}$, but no general proof was known. The Bloch--Kato conjecture, generalizing Milnor's to $\mathbb{Z}/l\mathbb{Z}$ coefficients for all primes $l$, remained completely open.

The challenge was to construct a cohomology theory for varieties over arbitrary fields that would be $\mathbb{A}^1$-invariant (invariant under homotopy with respect to the affine line) and would relate naturally to algebraic K-theory. Without such a theory, both the Milnor conjecture and the Bloch--Kato conjecture seemed out of reach.

An essential difficulty was that classical homotopy theory, built on topological spaces, has no direct analogue in algebraic geometry. The affine line $\mathbb{A}^1$ plays the role of the contractible interval, but this requires a suitable notion of ``weak equivalence'' in the algebraic category. Suslin had proposed using singular simplicial sets---replacing continuous maps from $\Delta^n$ with algebraic maps---but the resulting theory lacked key functorial properties. Voevodsky's breakthrough was to replace the topological model entirely with a sheaf-theoretic construction using the Nisnevich topology, which has enough points to detect all algebraic phenomena while being coarser than the \u00e9tale topology.

The Nisnevich topology is finer than the Zariski topology (it allows covers that are locally split in the \u00e9tale sense) but coarser than the \u00e9tale topology (it does not allow covers that are purely inseparable). Its key property is that it has enough points to classify vector bundles and other geometric objects, while still being ``algebraic'' enough to support a good theory of sheaves with transfers. This balance makes it the natural topology for motivic cohomology.

A key technical tool in Voevodsky's work is the \emph{cubical} approach to $\mathbb{A}^1$-homotopy theory, where instead of using simplicial schemes (which are topological in origin), one uses cubical schemes built from products of $\mathbb{A}^1$. This cubical structure makes $\mathbb{A}^1$-invariance automatic and simplifies the construction of Eilenberg--MacLane spaces and other homotopy-theoretic objects.

\section{How Was It Solved}

Voevodsky's proof of the Milnor conjecture proceeds in three major steps: constructing motivic cohomology, establishing the comparison with K-theory, and proving the isomorphism.

\textbf{Step 1: Motivic cohomology.} Working with Andrei Suslin and Eric Friedlander, Voevodsky defined motivic cohomology using the category of presheaves with transfers on smooth schemes. A presheaf with transfers is a contravariant functor from the category of finite correspondences (smooth schemes with algebraic correspondences as morphisms) to abelian groups. Motivic cohomology is then the sheaf cohomology in the Nisnevich topology on this category. The resulting groups are denoted:
\[
H^{p,q}(X, \mathbb{Z}) = \mathbb{H}^p(X, \mathbb{Z}(q))
\]
where $\mathbb{Z}(q)$ is the motivic complex, constructed inductively by $\mathbb{Z}(q) = \mathbb{Z}(q-1) \otimes \mathbb{G}_m[-1]$. This theory satisfies all expected properties: long exact sequences, Mayer--Vietoris, product structure, and---crucially---$\mathbb{A}^1$-invariance: $H^{p,q}(\mathbb{A}^1, \mathbb{Z}) = H^{p,q}(\operatorname{Spec} F, \mathbb{Z})$.

The definition of finite correspondences is central. A finite correspondence from $X$ to $Y$ is an algebraic cycle on $X \times Y$ that is finite and surjective over a component of $X$. Composition is given by intersection product. This category, denoted $\operatorname{Cor}(k)$, is much larger than the category of smooth schemes, and it is this enlargement that makes $\mathbb{A}^1$-invariance work: any morphism of schemes $\mathbb{A}^1 \times X \to Y$ gives a finite correspondence, and the sheaf condition in the Nisnevich topology ensures that presheaves with transfers are automatically $\mathbb{A}^1$-invariant.

\textbf{Step 2: Comparison with K-theory.} Voevodsky and Friedlander proved the fundamental comparison: for a field $F$,
\[
H^{n,n}(\operatorname{Spec} F, \mathbb{Z}) \cong K^M_n(F),
\]
identifying the diagonal motivic cohomology with Milnor K-theory. The key tool is the Atiyah--Hirzebruch spectral sequence:
\[
E_2^{p,q} = H^{p-q}(X, \mathbb{Z}(-q)) \Rightarrow K_{-p-q}(X).
\]

The comparison is established through a series of steps: first, showing that the $\mathbb{A}^1$-homotopy type of $\operatorname{Spec} F$ carries a canonical filtration whose graded pieces are motivic Eilenberg--MacLane spaces; then verifying that the motivic cohomology groups agree with the Milnor K-groups on these graded pieces via explicit computations with symbols and cup products.

\textbf{Step 3: The Milnor conjecture.} The isomorphism $K^M_n(F)/2 \cong H^n(F, \mathbb{Z}/2\mathbb{Z})$ follows by combining the comparison with K-theory with Deligne's theorem on the Weil conjectures. The norm residue map $K^M_n(F)/2 \to H^n(F, \mathbb{Z}/2\mathbb{Z})$ sends $\{a_1, \ldots, a_n\}$ to the cup product $(a_1) \cup \cdots \cup (a_n)$, and Voevodsky proves this map is an isomorphism using the motivic Steenrod algebra. The motivic Steenrod algebra is the algebra of cohomology operations on motivic cohomology, generated by the Steenrod squares $\operatorname{Sq}^i$ and Bockstein homomorphisms. Voevodsky shows that these operations satisfy the Adem relations and compute them explicitly in terms of the motivic filtration on algebraic cycles.

\textbf{Step 4: The Bloch--Kato conjecture.} Voevodsky, together with Markus Rost, extended the proof to all primes $l$, establishing the norm residue isomorphism:
\[
K^M_n(F)/l \cong H^n(F, \mu_l^{\otimes n}).
\]
The proof, completed in 2011 and running to several hundred pages, uses the motivic Steenrod algebra, the theory of symmetric powers, and the Rost--Voevodsky theorem on the motivic cohomology of fields. The proof of the Bloch--Kato conjecture is one of the important results in arithmetic geometry. A central ingredient is the \emph{degree-$n$ Kato conjecture}, which concerns the Galois cohomology of local fields and the structure of the cup product $H^1(F, \mu_l) \otimes \cdots \otimes H^1(F, \mu_l) \to H^n(F, \mu_l^{\otimes n})$.

\section{Impact of This Discovery}

The proof of the Bloch--Kato conjecture completed the program of ``higher class field theory''---the description of abelian extensions of a field in terms of its K-groups. It implies that the Galois cohomology of a field is completely determined by its K-theory, providing a bridge between algebraic K-theory and \u00e9tale cohomology that has applications throughout arithmetic geometry. The proof is one of the deepest results in modern algebraic geometry, combining motivic cohomology, the theory of symmetric powers, and the motivic Steenrod algebra in an essential way. The Bloch--Kato conjecture has implications for the structure of Galois groups of local and global fields, and it provides the theoretical foundation for the higher class field theory program initiated by Artin, Tate, and class field theorists in the early twentieth century.

Voevodsky's construction of motivic cohomology is now a standard tool. It is used for computing K-groups via spectral sequences, studying the arithmetic of varieties through the relationship between K-theory and \u00e9tale cohomology, and in birational geometry. The category of mixed motives, constructed by Voevodsky using triangulated categories of presheaves with transfers, provides a concrete realization of Grothendieck's vision.

The $\mathbb{A}^1$-homotopy theory that Voevodsky created---the homotopy category of schemes obtained by making $\mathbb{A}^1$ contractible and using the Nisnevich topology---has become an active area of research with applications to the cohomology of algebraic varieties, birational geometry, representation theory, and number theory. The stable motivic homotopy category $\mathcal{SH}(S)$ is a triangulated tensor category containing all known cohomology theories for algebraic varieties.

In his later years, Voevodsky turned to the foundations of mathematics, developing homotopy type theory (HoTT) and the univalence axiom $(A \simeq B) \simeq (A = B)$. He initiated the UniMath project for computer-verified mathematics, believing that increasing mathematical complexity required formal verification. His work on HoTT may ultimately have the greatest long-term impact.

The Bloch--Kato conjecture also implies the \emph{Quillen--Lichtenbaum conjecture}, which relates K-theory of schemes to \u00e9tale cohomology in high degrees. This had been an open conjecture for decades, and its resolution via Voevodsky's methods demonstrates the power of the motivic framework for solving deep problems in algebraic geometry and number theory.

The proof of the Bloch--Kato conjecture proceeds by induction on the degree $n$ and the number of primes dividing $n$. The base case $n=1$ is the classical Kummer sequence, which identifies $K^M_1(F)/l$ with $F^\times / (F^\times)^l \cong H^1(F, \mu_l)$. The inductive step uses the \emph{Norm Residue Isomorphism Theorem}, which states that the cup product
\[
H^1(F, \mu_l) \otimes \cdots \otimes H^1(F, \mu_l) \to H^n(F, \mu_l^{\otimes n})
\]
factors through $K^M_n(F)/l \to H^n(F, \mu_l^{\otimes n})$, and that this map is an isomorphism. The proof of the Norm Residue Isomorphism Theorem is the heart of Voevodsky's work, and it uses the motivic Steenrod algebra in an essential way.

The \emph{motivic Steenrod algebra} is the algebra of cohomology operations on motivic cohomology, generated by the Steenrod squares $\operatorname{Sq}^i$ and Bockstein homomorphisms. Voevodsky showed that the motivic Steenrod algebra satisfies the Adem relations and computed its structure explicitly. The key tool is the \emph{norm residue homomorphism}, which relates the cup product in Galois cohomology to the product in motivic cohomology, and which is shown to be an isomorphism using the motivic Steenrod algebra and the theory of symmetric powers.

\section{Modern Views and Improvements}

The motivic Adams spectral sequence, built from Voevodsky's framework, has produced new results on the classical homotopy groups of spheres, using the motivic filtration to simplify computations. The motivic Thom spectrum MGL, the algebraic analog of complex cobordism, provides the universal oriented cohomology theory in algebraic geometry.

The \emph{motivic Adams spectral sequence} has $E_2$-page $E_2^{p,q,s} = \text{Ext}^{p,q,s}_{\mathcal{A}}(\mathbb{Z}/2, \mathbb{Z}/2)$, where $\mathcal{A}$ is the motivic Steenrod algebra and $s$ denotes the stem. Voevodsky computed the motivic Adams spectral sequence in the ``2-primary stable stems'' and showed that it converges to the motivic stable homotopy groups of spheres. The key result is that the motivic cohomology of $\text{Spec}(k)$ is a polynomial algebraa $\mathbb{Z}/2[\tau, \rho]$, where $\tau$ is the motivic element of degree $(0,1)$ and $\rho$ is the element of degree $(-1,0)$. This computation is essential for the proof of the Milnor conjecture, as it provides the input for the motivic Adams spectral sequence.

The \emph{motivic Thom spectrum} MGL is defined as the Thom spectrum of the universal oriented vector bundle over the motivic Grassmannian. The cohomology of MGL is the algebraic cobordism $\Omega^*(X)$, which is the universal oriented Borel--Moore homology theory for smooth schemes. Panin and Gille proved that $\Omega^*(X)$ is isomorphic to the Lazard ring tensored with the Chow ring, and that the formal group law of MGL is the universal formal group law. This result provides a geometric interpretation of the algebraic cobordism and connects Voevodsky's work to the theory of oriented cohomology theories in algebraic geometry.

The geometric Langlands program, developed by Beilinson and Drinfeld, reformulates the Langlands correspondence as an equivalence of categories: $\operatorname{D-modules on Bun}_G \leftrightarrow \operatorname{Quasi-coherent sheaves on Loc}_{{}^L G}$. Voevodsky's work on motivic cohomology provides essential computational tools for this program, particularly through the motivic spectral sequence and the computation of motivic cohomology groups of classifying spaces.

The univalence axiom has become a cornerstone of homotopy type theory, with implications for type theory, computer science, and the foundations of mathematics. The HoTT book, resulting from a seminar at the Institute for Advanced Study, presents Voevodsky's ideas as a new foundation for mathematics compatible with computer verification.

The \emph{univalence axiom} states that the canonical map $(A = B) \to (A \simeq B)$ from identity types to equivalence types is itself an equivalence. This axiom has remarkable consequences: it implies function extensionality (if two functions agree on all inputs, they are equal) and the existence of higher inductive types (types with constructors that create both points and paths). The univalence axiom is consistent with Martin--L\u00f6f type theory, and it provides a computational interpretation of the identity type that is compatible with the homotopy theory of topological spaces. The HoTT book presents Voevodsky's univalence axiom as a new foundational principle for mathematics, alongside the Zermelo--Fraenkel axioms and the Peano axioms.

The \emph{motivic Adams--Novikov spectral sequence} is a variant of the classical Adams--Novikov spectral sequence in the motivic setting. The $E_2$-page of this spectral sequence is computed in terms of motivic cohomology groups, and it converges to the motivic stable homotopy groups of spheres. Voevodsky computed the motivic Adams--Novikov spectral sequence in the ``2-primary stable stems'' and showed that it collapses at $E_3$ in low degrees. The key input is the computation of the motivic cohomology of $\text{Spec}(k)$ as a polynomial algebra $\mathbb{Z}/2[\tau, \rho]$, which provides the $E_2$-page of the spectral sequence. The motivic Adams--Novikov spectral sequence has been used to compute new stable homotopy groups and to detect new elements in the stable homotopy groups of spheres.

Current research continues to develop the theory of motivic spectra, motivic stable homotopy groups, and the relationship between motivic cohomology and algebraic K-theory. The motivic Steenrod algebra remains an active tool for detecting classes in motivic cohomology. Recent work by Panin, Gille, and others on algebraic cobordism and oriented Borel--Moore homomorphisms builds directly on Voevodsky's constructions, and the motivic Adams--Novikov spectral sequence has been used to compute new stable homotopy groups in the motivic setting.

The \emph{algebraic K-theory of schemes} is closely related to motivic cohomology through the motivic spectral sequence:
\[
E_2^{p,q} = H^{p-q}(X, \mathbb{Z}(-q)) \Rightarrow K_{-p-q}(X),
\]
which converges from motivic cohomology groups to algebraic K-groups. Voevodsky's proof of the Bloch--Kato conjecture provides the key input for the convergence of this spectral sequence, and the resulting computation of K-groups has applications throughout arithmetic geometry. In particular, the K-theory of number fields controls the Galois cohomology via the Quillen--Lichtenbaum conjecture, which was proved as a consequence of Voevodsky's work.

The \emph{motivic stable homotopy category} $\mathcal{SH}(k)$ is the triangulated tensor category of motivic spectra over a base field $k$. It is constructed as the stabilization of the $\mathbb{A}^1$-homotopy category: a motivic spectrum is a sequence of motivic spaces $\Sigma^{\infty,n} X$ with structure maps $\Sigma^{1,1} X_n \to X_{n+1}$. The motivic stable homotopy groups $\pi_{p,q}(\mathbb{S})$ are the homotopy groups of the motivic sphere spectrum, and they generalize the classical stable homotopy groups of spheres. Voevodsky computed the motivic Adams spectral sequence for the 2-primary motivic stable homotopy groups of spheres, showing that the spectral sequence collapses at $E_3$ in low degrees and providing new information about the classical stable homotopy groups.

\begin{center}\rule{0.5\linewidth}{0.5pt}\end{center}

\section*{Summary}
\addcontentsline{toc}{section}{Summary}

Voevodsky developed motivic cohomology and the $\mathbb{A}^1$-homotopy theory of schemes, proving the Milnor conjecture and the full Bloch--Kato conjecture. His later work on homotopy type theory and the univalence axiom created a new foundation for mathematics with profound implications for computer-verified proofs.
